% TOP_PROBLEM: {"problemId": "problem.polynomial-diameter-bound-for-polytopes", "canonicalTitle": "Polynomial Diameter Bound for Polytopes", "releaseRank": 213, "primaryDomain": "cryptography_coding_information_optimization", "primaryDomainLabel": "Cryptography, coding, information & optimization", "displayStatus": "open", "releaseStatus": "open"}
% !TeX program = lualatex
% Definition and sources only; this file is not an LLM solution attempt.
\documentclass[11pt]{article}
\usepackage[margin=25mm]{geometry}
\usepackage{fontspec}
\setmainfont{Latin Modern Roman}
\usepackage{amsmath,amssymb,mathtools}
\usepackage{unicode-math}
\setmathfont{Latin Modern Math}
\usepackage{xurl,hyperref}
\hypersetup{colorlinks=true,urlcolor=blue}
\setcounter{secnumdepth}{0}
\setlength{\parindent}{0pt}
\setlength{\parskip}{5pt}
\setlength{\emergencystretch}{3em}
\begin{document}
\section*{\texorpdfstring{Polynomial Diameter Bound for Polytopes}{Polynomial Diameter Bound for Polytopes}}\label{213}
\subsection{Definitions and mathematical statement}
A pointed polyhedron is an intersection of finitely many closed affine half-spaces in ${\mathbb{R}}^d$ containing no affine line; assume its affine dimension is $d$. Its graph has vertices the zero-dimensional faces and edges joining pairs that share a bounded one-dimensional face. Graph diameter is the maximum shortest edge-path distance between vertices. The target is a single polynomial $p$ such that
\[
 \operatorname{diam}(P)\le p(n,d)
\]
for every pointed $d$-dimensional polyhedron with $n$ facets, including unbounded polyhedra. Coordinates and their bit lengths do not enter the bound. This asks for a polynomial diameter estimate, not the stronger historical linear bound nor a polynomial-time pivot rule.
\par\smallskip\textit{Formulation and concept references:} \href{https://drops.dagstuhl.de/entities/document/10.4230/LIPIcs.SoCG.2022.18}{[S1]}.

\subsection{Short English statement}
Can the graph distance between any two vertices of a pointed polyhedron always be bounded by a polynomial in its dimension and number of facets?
\subsection{Sources}
\begin{itemize}
\item[S1]Bonnet, Dadush, Grupel, Huiberts, and Livshyts, SoCG 2022.
\url{https://drops.dagstuhl.de/entities/document/10.4230/LIPIcs.SoCG.2022.18}.
\textit{Formulation source.}
\end{itemize}
\end{document}
